\section{The scale choice and the finite-stage induction}
\label{sec:iteration}
The construction uses one prescribed scale sequence. The quantitative source estimate supplies a fixed exponent $D_d$
for the particle-map invariant and a finite list of polynomial construction
guards. All choices in this section are made after those exponents have
been fixed.

\subsection{The base flow and the order of choices}
Choose an oriented orthonormal basis $(p,q,n)$ of $E$ and let the full
matrix $L_B$ vanish on $E^\perp$ and satisfy
\[
 L_Bq=p+X^{-2}n,\qquad L_Bp=L_Bn=0.
\]
It is trace free. Lemma~\ref{lem:potential}, with a fixed radial cutoff,
gives a compact, odd, reflection-equivariant solenoidal datum equal to
$L_Bx$ near zero. Its Sobolev and Gevrey bounds are uniform for $X\geq1$.
Lemma~\ref{lem:local} gives its actual Euler solution $u_B$ on a uniform
short interval. Its Gevrey bounds also persist on a uniform shorter
interval: with fixed base order $q_d> (d+1)/2+2$ and radius $r(t)$,
set
\[
 \mathcal X_r=\sum_{n\geq0}\frac{r^n}{(n!)^2}
                  \sum_{|I|=n}\|\partial_Iu_B\|_{H^{q_d}},
 \qquad
 \mathcal Y_r=\sum_{n\geq0}\frac{nr^n}{(n!)^2}
                  \sum_{|I|=n}\|\partial_Iu_B\|_{H^{q_d}}.
\]
At each finite cutoff the differentiated transport cancellation and
factorial product convolution give
$\mathcal X_r'\leq C_d\mathcal X_r^2+
(r'/r+C_d\mathcal X_r/r)\mathcal Y_r$.
Choose $r'$ a sufficiently negative fixed constant. A short-time
bootstrap bounds $\mathcal X_r$ and keeps $r>0$ uniformly in the
cutoff. Passing to the increasing sums proves the claim. The projected
Euler equation controls the time derivative and material acceleration
at a smaller radius. The integral flow equation then bounds the actual
particle map and its first two time derivatives in the same manner.
This proof uses the fixed legal-dimensional Sobolev embedding only.
Write $K_B$ for a uniform upper bound on its pressure Hessian.

First fix $d$, the packet base order and all finite inverse, coefficient,
flow and source exponents, including $D_d$. Fix $C_M,C_H$ large enough
for the global gradient and Hessian induction below. These fix the
endpoint constant and then its small threshold $\zeta$.
Choose $D_0$ sufficiently large for the first geometric comparison;
$D_0=1000$ suffices for the displayed exponent $60$.
All seed mean-cutoff and boundary costs are fixed powers of $X$.
Choose $D_k$ larger than the resulting seed polynomial exponent divided
by each of the finitely many frequency margins. Then choose the starting
index $J$ to pay all normal-stage exponent comparisons, and finally
choose $X$ large enough for their finite constants and total error budgets.
No choice is altered during the induction.

Put $h_{J-2}=1$, $x_{J-1}=X$, and $x_j=j^2x_{j-1}$ for $j\geq J$.
The auxiliary value $h_{J-2}$ records the uniform scale of the base
solution. Prescribe
\begin{align}
 h_{J-1}=X^{D_0},\quad k_{J-1}=X^{D_k},\quad
 \rho_{J-1}=X^{-D_0},\quad \delta_{J-1}=X^{-(D_0+10)},
 \label{eq:seed-scales}\\
 h_j=e^{x_{j-1}/j^5},\quad k_j=e^{x_{j-1}/j^2},\quad
 \delta_j=e^{-x_{j-1}/j^3},\quad
 \rho_j=e^{-x_{j-1}/j^{7/2}}\quad(j\geq J).
 \label{eq:normal-scales}
\end{align}
The seed frequency is fixed by its polynomial source bound in $X$;
the source-envelope lemma pays the prescribed normal frequencies.
Use the explicit physical mean-cutoff parameters
\[
 \ell_*=1,\qquad r_*=\frac{\rho_{J-1}}{4a_d}
                   =\frac{X^{-D_0}}{4a_d}.
\]
For large $X$, $r_*\le1/4$. The splitting ball of radius
$a_dr_*/\ell_*$ is exactly the seed support ball of radius
$\rho_{J-1}/4$. These parameters are fixed functions of the seed value
before the final choice of $X$; they are retained throughout the
iteration and are not replaced by the shrinking $\rho_j$.
Their reciprocals and all subsequent boundary penalties are bounded by
fixed powers of $X$. The cost multiplier in the full mean reserve is
bounded by the dimensional constant $1+c_d4^{-d}$, so it causes no
circular change of the starting index or the source exponents.

\subsection{Time and geometric comparisons}
Let
\[
 W_j=\frac{3x_jx_{j-1}}{\sqrt{h_{j-1}}},\qquad
 S_{J-1}:=S_{\rm base}=2W_J=6J^2X^{2-D_0/2},\qquad t_{J-1}=0.
\]
At stage $j$, the already constructed parent fixes $a_j,\beta_j$.
Define
\begin{equation}\label{eq:actual-times}
 t_j=t_{j-1}+\frac{x_j}{\sqrt{\beta_ja_jh_{j-1}}},
 \qquad S_j=t_j+2W_{j+1}.
\end{equation}
When $a_j,\beta_jx_{j-1}^2\in[1/2,2]$,
$W_j/6\leq t_j-t_{j-1}\leq2W_j/3$.
For $j>J$ the exact identities
\[
 \log k_{j-1}=\frac{x_{j-1}}{(j-1)^4},\qquad
 \log h_{j-1}=\frac{x_{j-1}}{(j-1)^7}
\]
show that $W_{j+1}/W_j$ is arbitrarily small uniformly after $J$:
its logarithm is
\[
 2\log j+2\log(j+1)-\frac{x_{j-1}}{2j^5}
                    +\frac{x_{j-1}}{2(j-1)^7}.
\]
At $j=J$ the last term is instead $(D_0/2)\log X$.
The negative term dominates in both cases. Hence $S_j\leq S_{j-1}$,
and with the numerical majorant $\Theta_j=C(1+x_jx_{j-1})$ the extra scaled interval obeys
$2\sqrt{a_jh_{j-1}}W_{j+1}\leq\Theta_j^{-60}$.
Also $t_J\geq S_{\rm base}/12$, so at positive-history stages
$t_{j-1}^{-1}\leq12S_{\rm base}^{-1}\leq h_{j-1}$.

Use $E_*=k_{j-1}^{-1/4}$, $G_*=1+h_{j-2}$ and
$K_h=k_{j-1}^{D_d}$. Each part of \eqref{eq:geometric-error} satisfies
its smallness guard at the same scale choice. For the spatial term the
logarithm is at most
\[
 -\frac{x_{j-1}}{j^{7/2}}+
       \frac{c_dD_dx_{j-1}}{(j-1)^4}
       +\frac{Cx_{j-1}}{(j-1)^7}+C\log x_{j-1}
 \leq-\frac{x_{j-1}}{2j^{7/2}}.
\]
This estimate uses the inherited geometric cost, which excludes the new
$\delta_j^{-1}$; using the complete source list here would fail.
The old frequency error has logarithm
$-x_{j-1}/(4(j-1)^4)$. For $j\geq J+2$,
\[
 \log(\epsilon G_*^2)\leq
 -\frac{x_{j-1}}{2(j-1)^7}+
 \frac{2x_{j-1}}{(j-1)^2(j-2)^7}+C
 \leq-\frac{2x_{j-1}}{5(j-1)^7}+C.
\]
For $j=J$, $\epsilon\Theta^{61}\leq C(J)X^{122-D_0/2}$;
for $j=J+1$ the new exponential beats the old polynomial shear.
These cover the two exceptional stages. Multiplying by any of the
fixed powers of $x_{j-1},\Theta_j$ still gives the required smallness.
The same comparisons give
\begin{equation}\label{eq:shear-separation}
 h_j/h_{j-1}^2\longrightarrow\text{arbitrarily large uniformly},\qquad
 \frac{\sqrt{h_{j-1}}}{x_{j-1}x_j}\gg G_*+1,\qquad
 \rho_j<r_{\rm core}.
\end{equation}
Here $r_{\rm core}>0$ is the fixed physical seed core radius.
Each use of $\gg$ in this proof denotes a fixed finite inequality whose
constant was chosen before $J,X$.

\subsection{The summed budgets}
The relative size estimate for the selected packet bounds its history
and early-time gradient and pressure increments by
\[
 b_j=C_dh_jh_{j-1}K_h^{c_d}\Theta_j^{c_d}\epsilon^{-c_d}
                         e^{-b x_{j-1}}
       \leq e^{-b x_{j-1}/2}.
\]
This follows by dividing the logarithm of the prefactor by $x_{j-1}$;
it is bounded by $C_d/(j-1)^4+o(1)$, with the seed predecessor treated
as a polynomial in $X$. During the late interval the pressure sign
\eqref{eq:positive-flux}, with its ambient neighborhood version, and
$f_\delta'\geq-1$ give the upper-Hessian increment
$C_d\delta_jh_jh_{j-1}$. Its scale satisfies
\[
 \delta_jh_jh_{j-1}\leq
                  e^{-x_{j-1}/(2j^3)}.
\]
The exact packet remainders are bounded by $k_j^{-1/4}$.
Consequently valid nonnegative increment budgets are
\begin{equation}\label{eq:actual-increment-budgets}
 e_{p,j}=C_d(b_j+\delta_jh_jh_{j-1}+k_j^{-1/4}),\qquad
 e_{i,j}=C_d(b_j+k_j^{-1/4}).
\end{equation}
The second one applies at time zero. For every fixed $A>0$ the positive
sequence $x_{j-1}/j^A$ increases at least geometrically once $J$ is large.
Thus the sums of the displayed exponential errors, and of every fixed
negative power of $x_{j-1}$, tend to zero with $X\to\infty$.
Their finite partial sums therefore preserve both the desired upper
pressure bound and the fixed-cutoff reserve of
\eqref{eq:mean-reserve-increment}. This last budget includes the core
term $c_dB_cr_*^dS_{\rm base}$; it is not dropped.
Initially, with the explicit frozen cutoff above, that term is
\[
 O_d\!\left(J^2X^{2+(1/2-d)D_0}\right)\longrightarrow0.
\]
Indeed $B_c=O(X^{D_0})$, $r_*=O_d(X^{-D_0})$ and
$S_{\rm base}=6J^2X^{2-D_0/2}$. This proves the assertion directly for
$d\ge3$, without an unquantified later shrink of the physical cutoff.
The other reserve terms tend to zero with $S_{\rm base}$.
This leaves a fixed positive margin before summing the normal stages.
Choosing their total $\sum_j e_{i,j}\le1$ also gives
$B_{c,j}\le B_{c,\rm seed}+1=O_d(X^{D_0})$, so every later
boundary penalty remains within the same polynomial seed envelope.

\subsection{Seed and induction step}
Set $u_{J-2}=u_B$. Apply Lemma~\ref{lem:exact-packet-child} on
$[0,S_{\rm base}]$ with the seed scales, $m_0=p$, $v(0)=q$,
$\alpha=\delta_{J-1}h_{J-1}$ and zero mean boundary penalty.
For this seed solve take $r=0$, $B_c=0$ and a global base lower bound
$B_e$; the interior region is empty. Thus the required penalty is indeed
$\mathcal L=0$, and the seed inverse is coercive on the short interval
without a special core term.
At time zero $m\cdot M_Bv=1$ throughout its compact packet support.
The uniformly bounded base equations preserve $m\cdot M_Bv\geq1/2$
on the chosen short interval. Thus the seed upper-pressure cost is at
most $C_dX^{-10}+k_{J-1}^{-1/4}$.
At $t_0=0$ every forced high and mean grade has zero initial value;
only the primary and its actual tensor corrector remain. They are
supported in the fixed seed core. The interior gradient bound is
$B_c=C_Mh_{J-1}+2$; the exterior lower bound remains that of the base.
Equip this actual seed output with the already prescribed fixed cutoff
parameters and penalty $\mathcal L=C_{1,d}B_c$ for subsequent stages.
Their full reserve is strictly small by the explicit estimate above.

The seed expansion is, at the fixed origin,
\[
 M_{J-1}=M_B+h(t)q(t)p(t)^T+\mathcal E(t),\qquad
 |\mathcal E|\leq k_{J-1}^{-1/4},\qquad h(0)=h_{J-1}.
\]
The normalized frame is transported by the actual base gradient.
Its activation parameters are exactly $a_J=1$ and $\beta_J=X^{-2}$.
The actual lifted flow estimate gives its polynomial particle-map
bound at $k_{J-1}$. Thus the complete input state for stage $J$ consists
of two exact solutions, the actual frame and shear decomposition,
the global low bounds, the strict fixed-cutoff reserve and the polynomial
label bounds. The first activation time is zero, so this stage uses the forward
primary without a history interval.

Assume that state at stage $j$. Keep the actual parent on its available
long interval $[0,S_{j-1}]$ while forming the endpoint and geometric
source. Its actual scaled horizon is
$a_j(S_{j-1}-t_{j-1})/\epsilon_j$, bounded by the numerical majorant
$\Theta_j$. By $S_{j-1}-t_{j-1}=2W_j$, it contains the target and the
short retained extension, and the scalar comparison interval lies within these bounds.
The numerical majorant is not an assertion that coefficient paths have
been defined beyond their actual time interval. Only after these source
estimates and the target normalization are obtained do we restrict the
parent and coefficient fields to $[0,S_j]$ for the packet and correction.
Their coefficient constants remain those of the long parent.
All Duhamel and energy estimates may be integrated on the actual scaled
interval of length at most $\Theta_j$; no extension to the numerical
majorant is needed. The relative comparison with the target size is
used only through $S_j$, whose extension past the target is small,
and is not asserted on the whole longer parent interval. Along the
older central solution,
$\dot B=-B^2-H_{j-2}$, which gives the bounds used above on $B,\dot B$.
Lemma~\ref{lem:actual-endpoint} supplies the actual stationary history
for $j>J$, and the zero-time prescription supplies it for $j=J$.
The ambient comparison and source-envelope argument verify the
quantitative neighborhood, full transverse propagator and every
coefficient cost at the already prescribed $\rho_j,k_j$.
Take
\[
 \alpha_j=\frac{\delta_jh_j}{|m(t_j,0)||v(t_j,0)|}.
\]
The joined construction in Lemma~\ref{lem:exact-packet-child} then gives
an exact smooth odd reflection-equivariant solution $u_j$ on $[0,S_j]$.
Its full ambient tensor correctors, all means and the residual correction
are part of this same solution. The graph-flow identities give the
next actual particle-map invariant $K_{h,\rm next}\leq k_j^{D_d}$.

The global expansion gives
\[
 \|M_j-M_{j-1}\|_\infty\leq C_dh_j+k_j^{-1/4},\qquad
 \|H_j-H_{j-1}\|_\infty\leq C_dh_jh_{j-1}+k_j^{-1/4}.
\]
With \eqref{eq:shear-separation} these preserve
$\|M_j\|_\infty\leq C_Mh_j$ and
$\|H_j\|_\infty\leq C_Hh_jh_{j-1}$.
The one-sided budgets \eqref{eq:actual-increment-budgets} preserve the
pressure upper bound, the initial interior and exterior quadratic-form
bounds, and the full fixed-cutoff reserve. No locality of the pressure
is being used: its signed expansion and its error bound hold globally.

For $t\in[t_j,S_j]$, the new normalized pair $p_2,q_2$ satisfies
\[
 M_j=M_{j-1}+h(t)q_2p_2^T+\mathcal E_j,
 \qquad |\mathcal E_j|\leq k_j^{-1/4},\qquad h(t_j)=h_j.
\]
Its frame equations have older field $M_{j-1}$, by the exact ray and
velocity equations, and $h'/h=-(M_{j-1})_{p_2p_2}-(M_{j-1})_{q_2q_2}$.
Lemma~\ref{lem:central-renewal} bounds each relative change in $a_j$
by $C_dx_{j-1}^{-4}$, whose sum is arbitrarily small. Starting at one,
the product stays in $[1/2,2]$. The same lemma gives directly
$\beta_{j+1}x_j^2\in[1/2,2]$; that estimate is not multiplied across
stages. Its compression bound and \eqref{eq:shear-separation} give
$p_2\cdot M_{j-1}p_2<-1$, while $|\mathcal E_j|<1$.
Also $(C_Mh_{j-1}+1)/h_j<\zeta$. These prove the two actual activation
conditions for the next endpoint problem. The low bounds on the entire
history and $t_j^{-1}\leq h_j$ prove its history hypotheses.
These estimates establish every component of the induction state
for the constructed child.

\section{The limiting datum and endpoint}
\label{sec:endpoint}
The initial estimates of Section~\ref{sec:initial-summability}, applied
to the same joined finite families and the actual source envelope, give
\[
 \sum_{j\geq J}\|u_j(0)-u_{j-1}(0)\|_{H^m}<\infty
                         \qquad\text{for every fixed }m.
\]
The seed is one fixed smooth increment. Lemma~\ref{lem:common-support}
places all these initial velocities in one compact set. Sobolev embedding
at arbitrarily high fixed orders therefore proves that their single limit
$u_0$ is smooth, compactly supported and divergence free.

The target times increase to a finite positive $T_\infty\leq S_{\rm base}$.
At the central target the exact expansion gives
\[
 |\nabla u_j(t_j,0)|_{\rm F}
       \geq h_j-C_Mh_{j-1}-k_j^{-1/4}\longrightarrow\infty.
\]
Let $u$ be the unique maximal smooth Sobolev solution from $u_0$.
If its lifetime $T^*$ exceeded $T_\infty$, it would be smooth on a
closed interval containing all $t_j$, with bounded $H^{s+1}$ norm.
Apply Lemma~\ref{lem:stability} on the actual varying intervals
$[0,t_j]$ to $u_j$ and $u$. Their initial $H^s$ difference tends to
zero, so their gradients converge uniformly on those intervals.
The displayed divergence of the central gradients would contradict
boundedness of $\nabla u$ there. Thus $0<T^*\leq T_\infty<\infty$.
This argument does not assert equality of $T^*$ and $T_\infty$.
Lemma~\ref{lem:continuation} now supplies the exact Frobenius gradient
limsup and the once-counted antisymmetric-vorticity integral in $H_d$.
Lemma~\ref{lem:newton} fixes the literal pressure normalization, and
Lemma~\ref{lem:local} supplies every required Sobolev order on each
shorter closed time interval.

The quantitative joined-source estimate and the prescribed scale
sequence verify the hypotheses at every finite stage. The initial limit,
finite lifetime and continuation conclusions therefore prove
Theorem~\ref{target:Hd}, with no additional assumption on the dimension
or on a prospective solution.
